% ============================================================================
%  Паттерны Тьюринга в модели Грея–Скотта — параметрическая серия
%  Компиляция:  xelatex turing_report.tex   (дважды, для колонтитулов/ссылок)
%  Требуется:   XeLaTeX, пакеты fontspec, polyglossia, graphicx, booktabs,
%               xcolor, geometry, fancyhdr, caption; шрифты DejaVu.
%  В одной папке с этим файлом должны лежать PNG:
%    map_fk.png gallery.png map_diff2d.png hysteresis.png coexistence.png switch.png
% ============================================================================
\documentclass[11pt]{article}

\usepackage{fontspec}
\usepackage{polyglossia}
\setmainlanguage{russian}
\setotherlanguage{english}

% --- Шрифты (DejaVu: есть кириллица, математические символы, моноширинный) ---
\setmainfont{DejaVu Serif}[Scale=0.92]
\setsansfont{DejaVu Sans}[Scale=0.92]
\setmonofont{DejaVu Sans Mono}[Scale=0.86]

\usepackage{geometry}
\geometry{a4paper, left=18mm, right=18mm, top=18mm, bottom=20mm}

\usepackage{graphicx}
\usepackage{booktabs}
\usepackage{array}
\usepackage{amsmath}
\usepackage{xcolor}
\usepackage{fancyhdr}
\usepackage{caption}
\usepackage{titlesec}
\usepackage[hidelinks]{hyperref}

% --- Цвета в тон исходного отчёта ---
\definecolor{acmag}{HTML}{C0295E}
\definecolor{accyan}{HTML}{00A5C0}
\definecolor{acink}{HTML}{1A2028}
\definecolor{acgrey}{HTML}{5A6472}
\definecolor{acrule}{HTML}{C8D0DA}

\color{acink}

% --- Заголовки ---
\titleformat{\section}{\sffamily\bfseries\Large\color{acmag}}{\thesection.}{0.5em}{}
\titleformat{\subsection}{\sffamily\bfseries\large\color{accyan}}{}{0em}{}
\titlespacing{\section}{0pt}{10pt}{6pt}
\titlespacing{\subsection}{0pt}{8pt}{3pt}

% --- Подписи к рисункам ---
\captionsetup{font={small,it,color=acgrey}, labelfont={color=acgrey}, justification=centering}

% --- Колонтитул ---
\pagestyle{fancy}\fancyhf{}
\renewcommand{\headrulewidth}{0pt}
\fancyfoot[L]{\ttfamily\scriptsize\color{acgrey}Grey--Scott · параметрическая серия}
\fancyfoot[R]{\ttfamily\scriptsize\color{acgrey}стр. \thepage}

\newcommand{\hruleline}{\vspace{-2pt}{\color{acrule}\hrule height 0.6pt}\vspace{6pt}}

\begin{document}

% ======================= Титул =======================
\thispagestyle{fancy}
\vspace*{45mm}
\begin{center}
{\sffamily\bfseries\fontsize{23}{28}\selectfont\color{acmag}
Паттерны Тьюринга в модели Грея–Скотта}\\[5mm]
{\sffamily\large\color{acgrey}
Онтография пространства параметров:\\ от карты режимов до трикритической точки необратимости}\\[10mm]
{\sffamily\normalsize\color{acgrey}
Вычислительная серия на браузерном симуляторе (headless Chrome~131)}\\[55mm]
{\itshape\small\color{acgrey}
Реакция–диффузия · численный эксперимент · 128\,$\times$\,128, CPU, детерминированный посев}
\end{center}
\clearpage

% ======================= 1. Метод =======================
\section{Модель и метод}
\hruleline
Исследуется двухкомпонентная система реакции–диффузии Грея–Скотта на торе:
\begin{equation*}
\frac{\partial A}{\partial t}=D_A\nabla^2 A-AB^2+f(1-A),\qquad
\frac{\partial B}{\partial t}=D_B\nabla^2 B+AB^2-(k+f)B .
\end{equation*}
Здесь $A$~— субстрат («питание»), $B$~— автокатализатор; $f$ (feed)~— скорость подпитки,
$k$ (kill)~— скорость удаления. Шаг по времени $dt=0.6$, пятиточечный лапласиан, границы
тороидальные. Симуляция управляется напрямую через внедрённый в страницу программный слой:
форсирован CPU-режим (для чтения полей), скорость максимальная, начальный шум обнулён, посев
детерминирован (фиксированный ГПСЧ), так что при всех прогонах начальное условие идентично, а
различия~— чисто динамические.

\subsection{Классификатор}
После релаксации измерялся дрейф поля $B$~— среднее $|\Delta B|$ за шаг в нескольких окнах.
Стационарные конфигурации держат ${\sim}(0.5\text{--}1)\cdot10^{-6}$ и затухают; нестационарные~—
$2\cdot10^{-5}\ldots3\cdot10^{-4}$ и не затухают. Порог ${\approx}\,2\cdot10^{-5}$ разделяет их чисто.
Плотность интерфейса (доля соседних клеток по разные стороны уровня $B=0.2$) служит мерой тонкости
структуры; покрытие~— доля $B>0.2$.

Всего в серии: карта $8\times6$ по $(f,k)$, плоскость $5\times5$ по $(D_A,D_B)$, зонд порога
нуклеации, проверка робастности к посеву и двусторонний прогон гистерезиса по $k$ для десяти
значений~$f$.
\clearpage

% ======================= 2. Карта (f,k) =======================
\section{Карта режимов в плоскости $(f,k)$}
\hruleline
Топология классическая и распадается на четыре области. При низком $k$ / высоком $f$ поле сплошь
заливается компонентом $B$ (однородное состояние без структуры). При высоком $k$ / низком $f$
компонент $B$ вымирает. Между ними~— полоса паттернов: стационарные лабиринты и пятна на
нижне-правой диагонали и кластер нестационарных (хаос, саморепликация) при низком $f$ и среднем $k$.

\begin{figure}[h]\centering
\includegraphics[width=\linewidth,height=0.80\textheight,keepaspectratio]{map_fk.png}
\caption{Атлас $(f,k)$. $f$ растёт вниз, $k$ вправо. Рамка: синий~— стационар, красный~—
нестационар, зелёный~— однородное заполнение, серый~— вымирание, оранжевый~— пограничное.}
\end{figure}
\clearpage

\begin{figure}[h]\centering
\includegraphics[width=0.96\linewidth]{gallery.png}
\caption{Шесть характерных режимов в разрешении.}
\end{figure}
Два наблюдения имеют прямую социально-структурную интерпретацию. Во-первых, рост $f$ (питания) при
умеренном $k$ уводит систему в однородное состояние: избыток ресурса стирает подструктуры.
Во-вторых, рост $k$ (умирания) монотонно снижает покрытие~— лабиринт $\to$ пятна $\to$ вымирание:
«репрессивная» среда измельчает и атомизирует структуру вплоть до коллапса. Это дискретный аналог
засушливой последовательности «лабиринт $\to$ пятна $\to$ голая почва», управляемой стрессом.
\clearpage

% ======================= 3. Диффузия =======================
\section{Плоскость диффузии $(D_A,D_B)$}
\hruleline
Одномерный прогон по $D_B$ смешивает два эффекта; двумерный их разделяет. \textbf{Масштаб} структуры
задаёт величина диффузии: вдоль линии постоянного отношения $D_A/D_B=4$ удвоение обоих коэффициентов
роняет плотность интерфейса $0.347\to0.245$~— отношение ${\approx}\sqrt2$, то есть тьюринговская длина
волны $\lambda\sim\sqrt{D}$. \textbf{Существование} паттерна задаёт отношение: при $D_A/D_B\lesssim1.3$
паттерн гибнет (субстрат не растекается достаточно быстрее автокатализатора). Но и слишком быстрый
субстрат вреден: вся строка $D_A=2.0$~— однородная заливка.

\begin{figure}[h]\centering
\includegraphics[width=\linewidth,height=0.72\textheight,keepaspectratio]{map_diff2d.png}
\caption{Плоскость $(D_A,D_B)$ при $f=0.030$, $k=0.057$. $D_A$ вниз, $D_B$ вправо;
подпись $i$~— плотность интерфейса.}
\end{figure}
Вывод: отношение диффузий определяет «режим» (есть структура или нет), а их общий масштаб~—
характерный размер ячейки, независимо от режима $(f,k)$.
\clearpage

% ======================= 4. Нуклеация и робастность =======================
\section{Порог нуклеации и робастность посева}
\hruleline
Вопрос: не зависит ли карта от начальной раскладки зёрен и нет ли на границах скрытой бистабильности,
которую слабый посев проглядел?

\subsection{Порог нуклеации}
Одиночное зерно возрастающего размера ($2\ldots18$) на пяти пограничных точках (со стороны
однородного и со стороны вымирания) не зажигает паттерн ни при каком размере. Контрольная внутренняя
точка зажигается от самого мелкого зерна. Значит, метки «однородно»/«мёртво» на границах~— не скрытая
бистабильность: где паттерн возможен, он зажигается от крошечной затравки; посев никогда не был узким
местом.

\subsection{Робастность}
Шесть ячеек, по четыре независимых случайных раскладки каждая, дали идентичный класс во всех случаях
(100\%). Класс и масштаб определяются уравнением, а не начальным условием.

\begin{center}\small
\begin{tabular}{@{}lll@{}}
\toprule
ячейка & класс $\times$ 4 раскладки & стабильно \\ \midrule
$f0.020\ k0.052$ & нестационар $\times 4$ & да \\
$f0.026\ k0.052$ & нестационар $\times 4$ & да \\
$f0.030\ k0.057$ & стационар $\times 4$   & да \\
$f0.048\ k0.062$ & стационар $\times 4$   & да \\
$f0.040\ k0.057$ & однородно $\times 4$   & да \\
$f0.014\ k0.062$ & вымирание $\times 4$   & да \\ \bottomrule
\end{tabular}
\end{center}
\clearpage

% ======================= 5. Гистерезис =======================
\section{Гистерезис на границе однородное $\leftrightarrow$ паттерн}
\hruleline
Границу проходили в две стороны с переносом состояния (адиабатическое продолжение): ветку «вниз» вели
от готового паттерна, снижая $k$; ветку «вверх»~— от однородного состояния, увеличивая $k$, с малым
возмущением (амплитуда $0.02$). Расхождение веток~— окно бистабильности.

\begin{figure}[h]\centering
\includegraphics[width=\linewidth,height=0.52\textheight,keepaspectratio]{hysteresis.png}
\caption{Плотность интерфейса vs $k$. Синий~— ветка «вниз» от паттерна, малиновый~— «вверх» от
однородного. Вверху $f=0.030$ (ветки совпадают~— мягкий переход), внизу $f=0.040$ (широкое окно
бистабильности).}
\end{figure}
При $f=0.030$ ветки практически совпадают~— переход непрерывный и обратимый. При $f=0.040$ паттерн
держится вниз до $k\approx0.058$, тогда как возмущённое однородное состояние не ломается вплоть до
$k=0.063$: окно $k\in[0.058,0.063]$, где оба состояния устойчивы одновременно.

\begin{figure}[h]\centering
\includegraphics[width=0.82\linewidth]{coexistence.png}
\caption{Сосуществование при ОДНИХ параметрах $f=0.040$, $k=0.060$: из паттернового начального
условия живёт лабиринт ($i=0.16$), из однородного (даже подтолкнутого) остаётся заливка ($i=0.00$).}
\end{figure}
\clearpage

% ======================= 6. Трикритическая точка =======================
\section{Трикритическая точка: где мягкий переход становится жёстким}
\hruleline
Порядковый параметр субкритичности~— амплитуда интерфейса в самый момент схлопывания паттерна на ветке
«вниз». В мягком (суперкритическом) сценарии структура истончается до нуля непрерывно; в жёстком
(субкритическом)~— срывается с конечной амплитудой.

\begin{figure}[h]\centering
\includegraphics[width=\linewidth,height=0.50\textheight,keepaspectratio]{switch.png}
\caption{Амплитуда в точке схлопывания vs $f$. Скачок из ${\sim}0.02$ ($f=0.028$) в конечные значения
($0.06\text{--}0.16$) при $f\ge0.030$ отмечает переключение суб-\,$\leftrightarrow$\,суперкритичности.}
\end{figure}

\begin{center}\small
\begin{tabular}{@{}lll@{}}
\toprule
$f$ & $k_{\text{fold}}$ (нижняя складка) & амплитуда в точке схлопывания \\ \midrule
0.028 & 0.0520 & 0.021\ \ (мягкий) \\
0.030 & 0.0535 & 0.063 \\
0.032 & 0.0555 & 0.156 \\
0.034 & 0.0560 & 0.081 \\
0.036 & 0.0570 & 0.132 \\
0.038 & 0.0575 & 0.117 \\
0.040 & 0.0580 & 0.105 \\
0.044 & 0.0590 & 0.114 \\ \bottomrule
\end{tabular}
\end{center}

\vspace{4mm}
Точка, разделяющая суперкритический и субкритический тьюринговский срыв,~— трикритическая
(коразмерности~2)~— лежит при $f\approx0.029\text{--}0.030$ ($D_A=1.0$, $D_B=0.5$). Ниже неё разрушение
порядка непрерывно и обратимо; выше~— дискретно и гистерезисно. Это согласуется с тем, что
локализованные и саморазмножающиеся структуры Грея–Скотта живут именно в субкритической зоне, где
паттерн сосуществует с устойчивым однородным фоном.
\clearpage

% ======================= 7. Синтез =======================
\section{Синтез}
\hruleline
Пространство параметров $(f,k,D_A,D_B)$ оказывается структурированным объектом с собственной
топологией: четыре области (однородное, вымирание, стационарные и нестационарные паттерны),
разделённые границами разного характера. Ключевые результаты серии:

\begin{itemize}\itemsep2pt
\item $f$ (питание) и $k$ (умирание) задают тип режима; отношение $D_A/D_B$~— само существование
структуры; общий масштаб диффузии~— её характерный размер ($\lambda\sim\sqrt{D}$).
\item Класс конфигурации инвариантен к начальному условию (робастность 100\%); порог нуклеации не
скрывает бистабильности на границах.
\item Граница однородное\,$\leftrightarrow$\,паттерн бывает мягкой (обратимой) и жёсткой
(гистерезисной), и переключение между ними~— трикритическая точка на оси питания при
$f\approx0.029\text{--}0.030$.
\end{itemize}

\subsection{Содержательный итог}
Получается двухуровневая «этика порога». Перейти складку $k$ (умирание) и разрушить структуру можно в
любой точке. Но обратимость самого этого разрушения включается параметром благополучия $f$: ниже
трикритической точки порядок восстанавливается простым откатом условий назад; выше~— однородное
состояние линейно устойчиво, и упавшую структуру одним возвратом $k$ уже не поднять. «Порог
необратимости»~— это не точка на оси умирания, а точка на оси благополучия: богатая среда даёт более
сложные устойчивые структуры и одновременно делает их гибель необратимой. Это прямая, посчитанная
реализация необратимого критического перехода~— и, пожалуй, наименее очевидный вывод всей серии.

\vspace{6mm}\hruleline
{\small\itshape\color{acgrey}
Приложения (сырые данные): \texttt{results.jsonl}~— карта $(f,k)$; \texttt{diff2.jsonl}~— плоскость
диффузии; \texttt{nucleation.jsonl}, \texttt{robust.jsonl}~— нуклеация и робастность;
\texttt{hysteresis.jsonl}, \texttt{hyst2.jsonl}~— ветки гистерезиса.}

\end{document}
